% 精度充足性
\section{实验四：精度充足性——什么时候L0就够了？}\label{sec:fidelity}

\textbf{论证目标}：验证梯度配置的实用价值——对于哪些设计点，低精度已经足够？对于哪些，需要升级？

\textbf{论证逻辑}：
梯度配置的价值不在于"我们有多个精度"，而在于"不同场景用不同精度的判断是一致的，
且在已知安全裕度的情况下可以自信地使用低精度"。

\textbf{实验设计}：
在$N=2, 6, 12, 18, 24$的Dragonfly配置上，分别用$(0,0,0)$和$(1,1,1)$精度求解，
比较$t^*$和feasibility判定。核心指标：L0的$t^*$相比L1是更保守（更安全）还是更乐观（有风险）？

\textbf{预期发现}：
\begin{itemize}
    \item L0始终更保守（更高的$t^*$，更低的$B_{\max}$）——二分带宽下的下界比Valiant分流更松
    \item 这意味着L0给出的"不可行"判断是可靠的（通过L0淘汰的，L1一定也淘汰），
    但L0的"可行"判断需要L1验证（L0可行的，L1可能仍不可行）
    \item 当L0的slack大于某个阈值时（如100×），升级到L1不改变判定，L0精度充足
\end{itemize}

\textbf{框架的使用模式}：先跑L0 $\to$ 观察slack $\to$ 仅当slack低于阈值时升级到L1。
精度升级成为数据驱动的决策，而非方法论绑定。
